%! Piecewise \& Step Functions — Unit 1, Lesson 1.4 — Lesson Plan
% GRADUAL-RELEASE plan (warm-up / guided notes & practice / debrief / close & start homework) on
% the stats/SAAR two-section shape, adopted for 1.4 on 2026-09-07. THERE IS NO GROUP ACTIVITY —
% the guided notes carry the release: I do (two sections) -> we do (guided practice). THERE IS NO
% INDEPENDENT PRACTICE SET IN THE NOTES — the homework is the individual practice, and the period
% ends with students starting it. No tiers, no exit ticket, no extension box. Teacher-facing,
% 10pt; every student component is 12pt. Timing 5 / 35 / 10 / 10 (SAAR, user decision
% 2026-09-06, the same 60-minute period).
\documentclass[10pt]{article}
\usepackage{algebra2-article}
\usepackage{algebra2-boxes}
\usepackage{graphicx}   % -article does NOT load graphicx; needed for thumbnails/screenshots
\graphicspath{{images/}}
\newcommand{\TallMath}[1]{$\displaystyle #1\rule[-1.4em]{0pt}{3.2em}$}
\newcommand{\CourseName}{Algebra 2}
\newcommand{\MeetingLength}{60 minutes}
\newcommand{\UnitNumberName}{Unit 1: Linear Functions \quad}
\newcommand{\LessonNumberName}{Lesson 1.4: Piecewise \& Step Functions}

\begin{document}
\begin{center}
    {\Large\bfseries \CourseName \\
    {\small\bfseries \UnitNumberName \LessonNumberName}}
\end{center}

\begin{tcolorbox}[colback=forestbg, colframe=forest, boxrule=0.9pt, arc=2mm,
    left=2mm, right=2mm, top=1.5mm, bottom=1.5mm]
    \textbf{Primary Objective:} Students read a \textbf{piecewise-defined function} as \emph{one}
    function built from different rules on different pieces of its domain. They \textbf{evaluate}
    it by first choosing the piece the input belongs to --- and at a \textbf{boundary}, the piece
    whose inequality \emph{includes} the input. They write the absolute value function as a
    two-piece linear rule, so Lesson~1.3's V is the first piecewise function they already know.
    They read a pre-drawn piecewise graph using \textbf{open} and \textbf{closed} endpoints, decide
    whether the pieces \emph{meet} at a boundary or \emph{jump} (a \textbf{discontinuity}), name
    increasing, decreasing, and \textbf{constant} intervals, and evaluate the
    \textbf{greatest-integer (step) function} $\lfloor x\rfloor$ --- rounding \emph{down},
    including for negatives.
    \\[2pt]
    \textbf{Standards (2023 VA SOL):} \textbf{A2.F.2} --- investigate and analyze the
    characteristics of \emph{piecewise-defined functions} algebraically and graphically: domain,
    range, zeros, and intercepts, \textbf{including graphs with discontinuities}
    (\textbf{A2.F.2a}); compare and contrast characteristics of piecewise-defined functions with
    the other families (\textbf{A2.F.2b}); increasing, decreasing, or \textbf{constant} intervals
    (\textbf{A2.F.2c}); and, for any $x$ in the domain, determine $f(x)$ from a graph or equation
    and explain its meaning in context (\textbf{A2.F.2f}). Absolute value (1.3) and the
    greatest-integer function are the two named entry examples of the piecewise family.
    \\[2pt]
    \textbf{Lesson model:} \emph{Gradual release --- I do, we do, you do --- all of it inside the
    guided notes.} There is no group activity. The notes name the vocabulary as they teach it and
    deliver the instruction in \textbf{two} sections on one context (a streaming service, free for
    three months); the Guided Practice is worked together with students holding the pen;
    \textbf{the homework is the individual practice}, and students start it in class while you
    circulate. The whole-class debrief is \emph{spoken} --- there is no practice set and no
    objective box on the notes page; the cover carries the targets.
\end{tcolorbox}

\renewcommand{\arraystretch}{1.4}
\boxguard
\begin{skillbox}[Priority Ideas \& Skills]{goldbox}
\begin{tabularx}{\linewidth}{@{}X|X@{}}
\textbf{Priority Ideas \& Skills} & \textbf{Key Understandings (the ``why'')} \\[2pt]
\hline
\vspace{-6pt}
\begin{itemize}[leftmargin=1.1em, nosep, after=\vspace{2pt}]
  \item Read brace notation: one function, several rules, each with the inputs it is for.
  \item \textbf{Evaluate} by choosing the piece first --- and at a boundary, the piece whose
        inequality includes the input.
  \item Write $|x|$ and $a|x-h|+k$ as a two-piece linear rule.
  \item Read a piecewise graph: open vs.\ closed dots; \emph{meets} or \emph{jumps}.
  \item Name increasing, decreasing, and \textbf{constant} intervals; write a rule from a graph.
  \item Evaluate $\lfloor x\rfloor$ (round \emph{down}) and read the staircase.
\end{itemize}
&
\vspace{-6pt}
\begin{itemize}[leftmargin=1.1em, nosep, after=\vspace{2pt}]
  \item A function can follow one rule here and another there and still be \emph{one}
        input-to-output machine --- the condition is a filter that sends each input to exactly
        one rule.
  \item \textbf{Target misconception:} that a boundary input belongs to ``both'' pieces, or to
        whichever piece the story makes sound right. Exactly one piece owns it --- the one whose
        inequality includes it --- and the closed dot shows which. The notes' gold box ($P(3)=0$)
        and Guided Practice (d) both test it.
  \item Pieces \emph{meet} only when both rules agree at the boundary; a jump is allowed, and it
        is what A2.F.2a means by a discontinuity.
  \item A step function is \emph{constant} on each interval, then jumps; ``rounding down'' is a
        step function, and its steps obey the same boundary rule.
\end{itemize}
\end{tabularx}
\end{skillbox}

\boxguard
\begin{skillbox}[{Vocabulary, Concepts, \& Theorems --- taught directly in the guided notes}]{sky}
\renewcommand{\arraystretch}{1.3}
\begin{tabularx}{\linewidth}{@{}l X@{}}
\textbf{Piecewise-defined function} & One function given by different rules on different parts of its domain, written with a brace. \\
\textbf{Boundary (breakpoint)} & An $x$-value where the rule changes; the piece whose inequality includes it ($\le$, $\ge$) owns it. \\
\textbf{Open / closed endpoint} & Open $\circ$ excludes the point ($<$, $>$); closed $\bullet$ includes it ($\le$, $\ge$). \\
\textbf{Continuous at a boundary} & Both rules give the \emph{same} value there, so the pieces meet with no gap. \\
\textbf{Discontinuity (jump)} & The two rules disagree at the boundary, so the graph jumps. \\
\textbf{Greatest-integer (step) function} & $\lfloor x\rfloor=$ the greatest integer $\le x$ --- rounds \emph{down}: $\lfloor 3.7\rfloor=3$, $\lfloor-2.5\rfloor=-3$; constant on each $[n,n+1)$, then jumps --- a staircase. \\
\end{tabularx}
\end{skillbox}

\begin{fixedskillbox}[Lesson at a Glance --- \MeetingLength]{forestbg}
\renewcommand{\arraystretch}{1.25}
\begin{tabularx}{\linewidth}{@{}l c X X@{}}
\textbf{Phase} & \textbf{Min} & \textbf{Students} & \textbf{Teacher} \\
\hline
Warm-Up & 5 & Three prior skills, alone, then check with a neighbour & Circulate; debrief item~3
  aloud (``$x=1$ fails $x<1$ but satisfies $x\le1$ --- so which rule gets it?'') and leave it
  hanging \\
Guided Notes \& Practice & 35 & Fill the vocabulary box and notes 1--2 as each term is named;
  then the Guided Practice together, pen in their hand & \textbf{I do} the hook and notes 1--2
  (20) $\to$ \textbf{we do} the Guided Practice box (15): $p(x)$, then change one piece and
  watch the boundary \\
Debrief & 10 & Pens down, papers up --- answer aloud, nothing to fill in & Put $C$ and $P$ back
  on the board; take the ``$P(3)=12$'' reflex and the round-toward-zero error \\
Close \& Start Homework & 10 & \textbf{Start the homework in class}, alone and silent & Launch
  item~1 aloud, then circulate for your formative read; preview Lesson~1.5; the rest is due the
  first class after two study halls \\
\end{tabularx}
\end{fixedskillbox}

\begin{fixedskillbox}[Warm-Up --- Activate Prior Knowledge (5 min)]{forestbg}
\begin{minipage}[t]{0.48\linewidth}
    \textbf{The three items and what each seeds}
    \begin{itemize}
        \item \textbf{1.} $g(x)=2x-1$ at $-2,0,3$ and $h(x)=5-x$ at $2,-4$. \textbf{Seeds:} the
              mechanical skill every \emph{piece} needs --- a piecewise function is several of
              these, each with a condition.
        \item \textbf{2.} The V as two lines, $y=x$ for $x\ge0$ and $y=-x$ for $x<0$, with the
              check $|-6|=-(-6)$. \textbf{Seeds:} the whole bridge from Lesson~1.3 --- section~1's
              second half writes it with a brace.
        \item \textbf{3.} Open vs.\ closed dots, $<$ vs.\ $\le$, ending on \emph{if one rule is
              for $x<1$ and another for $x\ge1$, which rule does $x=1$ use?} \textbf{Seeds:} the
              boundary rule, which section~1 states and section~2's crux tests.
    \end{itemize}
\end{minipage}
\hfill
\begin{minipage}[t]{0.48\linewidth}
    \textbf{Running it}
    \begin{itemize}
        \item \textbf{Debrief item~3 aloud, briefly}, and land only this: \emph{$x=1$ fails $x<1$
              and satisfies $x\le1$ --- hold on to which rule that gives it.} Do not state the
              boundary rule yet; the hook's month-3 vote is worth more than the ten seconds.
        \item Item~2's check, $|-6|=-(-6)=6$, is the one that convinces students the ``$-x$''
              piece is not negative. Do not rush past it.
        \item If it runs long, skip item~2's check and keep item~3 whole.
        \item The page is 12pt like the rest of the packet and fits one side. If you add to it,
              check that it still does.
    \end{itemize}
\end{minipage}
\end{fixedskillbox}

\boxguard
\begin{skillbox}[Guided Notes \& Practice --- I do, then we do (35 min)]{forestbg}
\textbf{This is the whole lesson.} There is no group activity, and there is no independent
practice set on this page: \textbf{the homework is the individual practice}, started in class.
\begin{multicols}{2}
\textbf{Hook (3 min, page 1).} ``A streaming service is free for the first 3 months, then \$12 a
month. Total paid after 2 months? After 5?'' (\$0, \$24.) ``Could one straight line describe the
total?'' (No --- flat, then rising.) Then the vote that carries the lesson: \textbf{``Month 3 ---
the last free month, or the first paid one?''} Take a hand count on \emph{free} vs.\ \emph{\$12},
hear one reason for each, and \emph{do not settle it}. Section~1 settles it, and section~2's gold
box makes the wrong vote cost something. Fill the vocabulary rows \emph{as you go}, never in
advance.

\textbf{I do (about 17 min).} Two sections, each in two moves.

\emph{Section 1 (page 2).} Read the brace for $C(m)$ --- each rule comes with the inputs it is
\emph{for}. Evaluate by \emph{choosing the piece first}: $C(2)=0$, $C(5)=24$ in the first
\texttt{work} block, then $C(3)$ --- $3$ satisfies $0\le m\le3$, the $\le$ includes it, so month~3
is still free: the vote is settled. State the \textbf{boundary rule} and its last blank aloud:
\emph{the story does not decide it; the inequality does.} Five minutes. Its second half, \emph{The
V Is Two Pieces}, writes $|x|$ with a brace, checks $|-3|$, and rewrites the trail
$H(x)=2|x-3|+1$ as $2x-5$ / $7-2x$ in two side-by-side \texttt{work} blocks; both give $H(3)=1$,
so the pieces \emph{meet} --- the first ``meets'' example, and section~2's contrast needs it.
Four minutes.

\textbf{Section 2 (page 3) is the lesson.} Its first half is the two-column table on the
\emph{same} story: total paid $C(m)$ meets at $3$ (both rules give $0$; constant on $[0,3]$,
increasing after) beside price-this-month $P(m)$, which jumps ($0$ and $12$; range $\{0,12\}$).
Land the one test: \emph{both rules agree at the boundary, or they do not.} Three minutes. Its
second half is the crux --- the gold box: a classmate reads $P(3)=12$ ``because the price is \$12
after the free months.'' But $3$ satisfies $0\le m\le3$, so $P(3)=0$, and the closed dot at
$(3,0)$ says so; the open dot at $(3,12)$ is not a point of the graph. \emph{Exactly one piece
owns the boundary.} Say it, write it, and point at the two dots while you do. Then \emph{Step
Functions: Rounding Down} --- $P(m)$ is constant, then jumps; $\lfloor x\rfloor$ rounds down,
$\lfloor-2.5\rfloor=-3$ (down means more negative); the staircase is constant $n$ on $[n,n+1)$,
closed left, open right --- the boundary rule again, one step at a time. Five minutes.

\columnbreak
\textbf{We do (about 15 min, page 4).} The Guided Practice box: $p(x)=-2x$ for $x\le-1$, $x+3$ for
$x>-1$, with its graph pre-drawn. Students hold the pen; you hold the questions. \textbf{(a)}
evaluate at $-3$, $0$, $2$, then $-1$ --- \emph{which piece owns it, and how do you know?}
\textbf{(b)} meets or jumps at $-1$ (both give $2$: meets). \textbf{(c)} decreasing on
$(-\infty,-1]$, increasing on $[-1,\infty)$, minimum $2$, read off the graph. \textbf{(d) is the
crux transferred and the part to protect:} a classmate says ``both rules give $2$, so it does not
matter which piece owns $-1$'' --- change the second piece to $x+4$. The first piece still owns
$-1$ ($p(-1)=2$), the second now gives $3$ just past it, the graph jumps, and the closed dot is at
$(-1,2)$: \emph{ownership is decided by the inequality, not by the rules happening to agree.}
\textbf{(e)} round down --- $\lfloor2.9\rfloor$, $\lfloor-1.2\rfloor$ (the trap: $-2$),
$\lfloor4\rfloor$, the interval $[4,5)$, and \emph{does that step own $5$?} (No: $[5,6)$ does,
$\lfloor5\rfloor=5$.)

\textbf{The pen must actually change hands.} Guided Practice is the only place today where you
can watch a student decide, so do not narrate it. Put part (a)'s $p(-1)$ on them cold, take a
hand count on \emph{meets or jumps} in part (d) before anyone speaks, and make the room defend
both answers.

Circulate during Guided Practice and demand the reason, not the word:
\begin{itemize}[leftmargin=1.1em, nosep]
  \item \textbf{Part (a):} ``Does $-1$ satisfy $x\le-1$? Does it satisfy $x>-1$? So which rule
        gets it?''
  \item \textbf{Part (d), the crux transferred:} ``You changed the second rule. Did you change
        the \emph{inequality}? Then who owns $-1$?''
  \item \textbf{Part (e):} ``$-1.2$ --- which integer is just \emph{below} it on the number line?''
\end{itemize}
While circulating, pick two papers to put up in the debrief --- one part (d) with the closed dot
at $(-1,2)$ and the open dot at $(-1,3)$, and one that moved $p(-1)$ to $3$.
\textbf{Your formative read comes twice}: here, and again in the last ten minutes as students
start the homework.
\end{multicols}
\end{skillbox}

\boxguard
\begin{skillbox}[Debrief --- whole class, spoken (10 min)]{forestbg}
Pens down, papers up. \textbf{This debrief is spoken} --- there is no box at the end of the
notes to fill in, so it lives entirely on the board and in the room.
\begin{multicols}{2}
\begin{enumerate}[leftmargin=1.3em, nosep, itemsep=3pt]
  \item \textbf{Put $C(m)$ and $P(m)$ back up side by side} with their dots, and make a student
        say which meets and which jumps --- and point at the two dots at $m=3$ that earned it.
        Then state the test plainly: \emph{same value at the boundary, they meet; different, it
        jumps.}
  \item Put up the two part (d)s you picked. The one that moved $p(-1)$ to $3$ read the
        \emph{open} dot; make the room say which inequality includes $-1$. That is the moment
        ``exactly one piece owns the boundary'' stops being a slogan --- point back at $P(3)=0$.
  \item Take the round-toward-zero error: ask who wrote $\lfloor-1.2\rfloor=-1$, and put $-1.2$
        on a number line --- \emph{down} means more negative. Then $\lfloor5\rfloor=5$: the step
        $[4,5)$ stops \emph{without} $5$, the next one starts \emph{with} it.
  \item Take the ``every piecewise graph is discontinuous'' reflex last: the trail $H(x)$ and
        $p(x)$ both meet; pieces \emph{can} agree.
\end{enumerate}
\columnbreak
\textbf{Two things to demand aloud.} First, \textbf{make a student say the month-3 sentence
without the notes} --- $3$ satisfies $0\le m\le3$, so the first piece owns it, so $P(3)=0$, and
the closed dot is at $(3,0)$. Second, \textbf{make them say what a closed dot means in words} ---
this endpoint is \emph{in} --- because homework~3 asks them to read one and then write the rule
behind it.

\textbf{If the debrief runs short}, cut item~4 and item~3 --- item~2 is the only one that cannot
be cut. \textbf{Do not let the debrief eat the last ten minutes} --- the homework start is not
optional padding, it is the individual practice.
\end{multicols}
\end{skillbox}

\boxguard
\begin{skillbox}[Homework --- scored, started in class, due the first class after two study halls]{goldbox}
Six items spanning A2.F.2a/b/c/f, two pages, no extension: \textbf{1} evaluating a two-piece rule
at four inputs including the boundary, with ``which piece owns $-1$, and why''; \textbf{2} the
deliberate \emph{meets-or-jumps} pair, one number apart ($g$ jumps, $3$ vs.\ $4$; $k$ meets), with
the ``what decides it'' sentence; \textbf{3} \emph{reading a graph with a jump and then writing
its rule} --- the only item that runs backwards, and the best single predictor of whether the
closed-dot reading stuck; \textbf{4} the staircase, the rounding traps $\lfloor-3.1\rfloor$ and
$\lfloor-0.5\rfloor$, and the boundary question ``when is $\lfloor x\rfloor=x$?''; \textbf{5} the
phone-plan model with a \texttt{work} block, an interpret-the-answer follow-up, and a surcharge
that opens a jump; \textbf{6} an SOL-style multiple-choice item evaluating at a boundary, the
formative check.

\textbf{Start it in the last ten minutes of the period, in class and alone.} Item~1's last blank
is the boundary check and item~3's $f(2)$ is where the misconception shows; the Close box says how
to sort what you see. \textbf{Item~6 carries the check when you score it:} (A) is the strict-$<$
piece (the target misconception), (C) is arithmetic, (D) thinks a boundary belongs to no one.

\textbf{Two pages is a ceiling at 12pt}, so the old extension items (a three-piece rule from a
description; a tax bracket that is continuous at \$10{,}000) were cut with this shape; use them
as debrief cold-calls if time allows. \textbf{DeltaMath override:} evaluating piecewise functions
and the greatest-integer function are well covered there --- swap in a DeltaMath set if you
prefer auto-graded practice; that replaces this page and strikes its score cell on the cover.
Never assign both; the graph-reading and continuity items are thinner there, so keep items~2
and~3 if you do. \textbf{Preview:} Lesson~1.5, \emph{Linear Regression} --- we leave exact rules
behind and fit one line to \emph{scattered} data.
\end{skillbox}

\boxguard
\begin{skillbox}[Watch For (while circulating)]{redbox}
\begin{itemize}[leftmargin=1.2em, nosep]
  \item \textbf{Giving a boundary input to the strict-$<$ piece, or to ``both''} (notes~1,
        GP~(a) and~(d), HW~1 and~6) --- the lesson's target misconception. Send them to the
        inequality, not to the story. Probe: \emph{``Does it satisfy $x\le2$? Does it satisfy
        $x>2$? So which rule gets it?''}
  \item \textbf{Reading the open dot as the value at a breakpoint} (notes~2's gold box,
        GP~(d), HW~3). Probe: \emph{``Two dots at that $x$. Which one is filled in?''}
  \item Calling every piecewise graph ``discontinuous'' --- pieces \emph{can} meet; check whether
        the two rules agree there (notes~1's trail, GP~(b), HW~2).
  \item Greatest-integer sign error: rounding \emph{toward zero}, so $\lfloor-1.2\rfloor=-1$
        (notes~2, GP~(e), HW~4). Probe: \emph{``Which integer is just below it?''}
  \item Saying a step function is ``increasing'' on a step --- it is \textbf{constant} there,
        then jumps (notes~2).
  \item In HW~5(c), not seeing that a surcharge opens a gap: the pieces meet only if both rules
        give \$30 at $5$~GB.
\end{itemize}
Cold-call prompts: ``Which piece owns this input --- how do you know?'' ``Open or closed here, and
what inequality does that show?'' ``Do the two rules agree at the boundary?'' ``Round down ---
which integer is just below?''
\end{skillbox}

\boxguard
\begin{skillbox}[Close \& Start Homework (10 min)]{goldbox}
\textbf{Students start the homework here, in class, alone and silent --- this is the individual
practice, and the first minutes of it are your best formative read.} Hand the launch first ---
six items on paper, scored, due the first class after two study halls (announce the date in class
and on TurtleNet) --- and work item~1's $f(-3)$ aloud in thirty seconds so nobody stalls on the
page. Then circulate on item~1's last blank and item~3's $f(2)$, sorting into three piles: chose
the piece by its inequality (got it); chose the strict piece or ``both'' (the target misconception,
still alive); read the open dot (send them to $P(3)$). Close on what changed: \emph{``You walked
in with functions that follow one rule. You walk out able to read a function that follows several
--- pick the piece, mind the boundary, read the dots --- and to see Lesson~1.3's V as the first
piecewise function you already knew.''}
\end{skillbox}

% ── Teacher notes — one per component, in packet order (three: no activity, no practice set) ──
\begin{teachernote}[Warm-Up]
Item~1 is the mechanical core: each piece of a piecewise function is a rule you evaluate, so
plugging in has to be automatic --- including into $5-x$ with a negative input. Item~2 is the
whole bridge from Lesson~1.3: the V is literally two linear rules joined at the vertex, and the
check $|-6|=-(-6)=6$ is the one that convinces students the ``$-x$'' piece is not negative.
Item~3's last bullet is the make-or-break idea for today --- $x=1$ fails $x<1$ and satisfies
$x\le1$, so it belongs to \emph{exactly one} rule. Debrief it in one sentence and leave it
hanging; the hook's month-3 vote and section~1's boundary rule pay it off within five minutes. The
page is 12pt like every other student component and fits one side, blank and keyed.
\end{teachernote}

\begin{teachernote}[Guided Notes \& Practice]
\textbf{This one document is the lesson --- pace it 3 / 17 / 15, then 10 for the debrief, then 10
for the homework start.} Page~1 is the vocabulary rows and the hook; fill the rows \emph{as each
term is named}, never in advance, and take the month-3 vote without resolving it. Section~1 must
move: five minutes on the brace, $C(5)$, and the boundary rule, four on the V as two pieces.
\textbf{Do not cut ``both rules give $H(3)=1$, so they meet''} --- it is the first meets example,
and section~2's contrast needs it.

\textbf{Spend the time in section~2}, which carries the misconception. Fill the two-column table
with the room on the \emph{same} story --- total paid meets, price this month jumps --- and read
the one test off the two pictures before naming anything. Then the gold box: let someone say
``$P(3)=12$, the price is \$12 after the free months'' and let the room catch it. Do not resolve it
by repeating the rule; point at the two dots at $m=3$. Close on the staircase --- four floor values,
the ``down means more negative'' sentence, and $[n,n+1)$ closed left, open right, which is the
boundary rule one step at a time.

\textbf{Guided Practice is where the pen changes hands.} Part (a) is the boundary check, (b) the
meets test, (c) the graph read, \textbf{(d) the crux transferred} --- change the second piece to
$x+4$ and the rules stop agreeing, but the \emph{inequality} has not changed, so the first piece
still owns $-1$ --- and (e) rounding down with ``does the step own $5$?'' If time is short, cut
(c), never (d).

\textbf{There is no independent practice set on this page, and that is deliberate --- the
homework is the individual practice.} Guided Practice is the last thing on the page; the period
ends with students starting the homework in front of you, which is where your second formative
read comes from. Three \texttt{work} blocks in the notes ($C(5)$ and the two halves of the trail).

\textbf{Debrief (10 min), spoken.} $C$ and $P$ side by side, then the two part (d)s; item~2 is the
one move that cannot be cut. \textbf{Do not let the debrief eat the last ten minutes} --- the
homework start is not optional padding.
\end{teachernote}

\begin{teachernote}[Homework]
\textbf{Scored; due the first class after two study halls --- announce the date in class and on
TurtleNet.} The ten minutes at the end of the period are a \emph{start}, not a completion: put
students on 1 and 2 while you can still catch a wrong start, and let 3--6 go home. Item~1's last
blank is the boundary check --- ``the first, because $-1\le-1$'' is the whole idea in six words;
expect $f(-1)=-2$ from anyone who used the strict piece. Item~2 is the meets-or-jumps pair in the
wild, and the sentence should say ``same value at the boundary.'' Item~3 is the one that runs
backwards --- $f(2)=-2$ off the closed dot, then the rule $x$ on $[-2,2)$ and $-2$ on $[2,5]$;
a student who writes $f(2)=2$ read the open dot. Item~4 has the two rounding traps
($\lfloor-3.1\rfloor=-4$, $\lfloor-0.5\rfloor=-1$) and the boundary question ---
$\lfloor x\rfloor=x$ exactly at the integers, never above $x$. Item~5(c) is the concept check: a
\$10 surcharge makes the second piece $40+10(g-5)$, which gives \$40 at $5$~GB against the first
piece's \$30 --- a jump. \textbf{When you score the page, sort item~6 into three piles}: chose B
(got it); chose A (used the strict-$<$ piece --- the target misconception, still alive); chose C
or D (arithmetic, or thinks a boundary belongs to no one). If more than a third land outside B,
open Lesson~1.5 with a one-minute re-look at ``which inequality includes it'' rather than letting
it ride. One \texttt{work} block (item~5b). \textbf{DeltaMath override:} if you swap in a
DeltaMath set, it replaces this page and strikes its score cell on the cover --- never assign both.
\end{teachernote}

\end{document}
